% ID: FIS-URT-ELA-001
% Titolo: Due urti elastici successivi tra tre palle
% Disciplina: Fisica
% Area: Quantita di moto e urti
% Argomento: Urti elastici
% Sottoargomento: Urti frontali successivi in una dimensione
% Classe: Terza liceo scientifico
% Difficolta: 4
% Tipo: problema_numerico
% Risultato: m_C = 0,597 kg; v_B'' = -10,8 m/s
% Tempo_stimato: 12 min
% Competenze: conservazione della quantita di moto, conservazione dell'energia cinetica, modellizzazione di urti successivi
% Prerequisiti: quantita di moto, energia cinetica, urto elastico unidimensionale
% Tag: fisica, quantita_moto, urti_elastici, urti_successivi, conservazione_energia
% Fonte: rielaborazione_da_traccia_fornita_dal_docente
% Autore: Riccardo Carli
% Licenza: CC-BY-SA-4.0

\begin{esercizio}
Una pallina di massa \(m_A=45{,}0\,\mathrm{g}\) si muove lungo una retta con
velocita \(v_A=135\,\mathrm{km/h}\) e urta elasticamente una seconda palla di
massa \(m_B=145\,\mathrm{g}\), inizialmente ferma. La seconda palla prosegue e
urta elasticamente, sempre frontalmente, una terza palla inizialmente ferma.
Dopo il secondo urto la terza palla si muove con velocita
\(v_C'=25{,}0\,\mathrm{km/h}\).

Determina:
\begin{enumerate}
\item la massa \(m_C\) della terza palla;
\item la velocita della seconda palla immediatamente dopo il secondo urto.
\end{enumerate}
\end{esercizio}

\begin{soluzione}
Convertiamo anzitutto le velocita nel Sistema Internazionale:
\[
v_A=\frac{135}{3{,}6}=37{,}5\,\mathrm{m/s},
\qquad
v_C'=\frac{25{,}0}{3{,}6}=6{,}94\,\mathrm{m/s}.
\]
Le masse sono
\[
m_A=0{,}0450\,\mathrm{kg},
\qquad
m_B=0{,}145\,\mathrm{kg}.
\]

Nel primo urto elastico la palla \(B\) e inizialmente ferma. Per un urto
elastico unidimensionale la sua velocita dopo l'urto e
\[
v_B'=\frac{2m_A}{m_A+m_B}v_A.
\]
Pertanto
\[
v_B'=\frac{2\cdot0{,}0450}{0{,}0450+0{,}145}\cdot37{,}5
\simeq17{,}76\,\mathrm{m/s}.
\]

Questa e la velocita con cui \(B\) raggiunge la palla \(C\). Nel secondo urto,
ancora elastico e con \(C\) inizialmente ferma,
\[
v_C'=\frac{2m_B}{m_B+m_C}v_B'.
\]
Ricaviamo \(m_C\):
\[
m_C=\frac{2m_Bv_B'}{v_C'}-m_B.
\]
Sostituendo i dati:
\[
m_C=\frac{2\cdot0{,}145\cdot17{,}76}{6{,}94}-0{,}145
\simeq0{,}597\,\mathrm{kg}.
\]

Per ottenere la velocita di \(B\) dopo il secondo urto usiamo la proprieta
dell'urto elastico unidimensionale secondo cui la velocita relativa di
allontanamento e uguale alla velocita relativa di avvicinamento:
\[
v_C'-v_B''=v_B'-0.
\]
Da cui
\[
v_B''=v_C'-v_B'=6{,}94-17{,}76
\simeq-10{,}8\,\mathrm{m/s}.
\]
Il segno negativo indica che la palla \(B\) inverte il verso di moto dopo il
secondo urto.
\end{soluzione}
